\subsubsection{Several-Variable Differentiation}

Several-variable differentiation studies best linear approximations and their local geometry. We use the linear-map, operator-norm, determinant, and invertibility tools from Proposition \ref{prop:matrix-representation-action}, Definition \ref{def:operator-norm-rn-rm}, Theorem \ref{thm:determinant-basic-laws}, and the metric background from Example \ref{ex:coordinate-norm-metrics-rk}. The operator-norm topology of invertible linear maps is developed here because it drives local inversion.

\begin{definition} \label{def:several-variable-differentiation-core-source-todo}\label{def:total-derivative-rn-rm} \hlt{Total Derivative}\\
Let $E\subseteq\R^n$ be open, $\f:E\to\R^m$, $\x\in E$, and $\LL(\R^n,\R^m)$ denote space of linear maps (Definition \ref{def:linear-map}).\\
If $\exists A\in\LL(\R^n,\R^m)$ s.t. $\lim\limits_{\h\to\ze}\frac{\norm{\f(\x+\h)-\f(\x)-A\h}_2}{\norm{\h}_2}=0$, then $\f$ is \hlt{differentiable at $\x$}, and $A$ denoted $D\f(\x)$ or $\f'(\x)$.\\
Equivalently, $\f(\x+\h)-\f(\x)=D\f(\x)\h+\mathbf{r}(\h)$, and $\frac{\norm{\mathbf{r}(\h)}_2}{\norm{\h}_2}\to0$.\\
If this holds for every $\x\in E$, then $\f$ is differentiable on $E$.
\end{definition}

\begin{perspective} \label{per:total-derivative-linear-map-first} \hlt{Derivative First, Matrix Second}\\
The total derivative is a linear map first; the Jacobian is its matrix in chosen bases. This keeps coordinate formulas subordinate to the intrinsic estimate $\norm{A\h}_2\leq\normop{A}\norm{\h}_2$ from Definition \ref{def:operator-norm-rn-rm}.
\end{perspective}

\begin{figure}[H]
  \centering
  \begin{tikzpicture}[
  scale=0.88,
  >=Stealth,
  every node/.style={font=\scriptsize},
  axis/.style={->, gray!70, line width=0.45pt},
  vector/.style={->, black, line width=0.75pt},
  linear/.style={->, green!45!black, line width=0.85pt},
  error/.style={->, black!65, dashed, dash pattern=on 3pt off 2pt, line width=0.55pt},
  point/.style={circle, fill=black, inner sep=1.15pt},
  approxpoint/.style={circle, fill=green!45!black, inner sep=1.15pt}
]
  \begin{scope}
    \draw[axis] (-0.25,0) -- (3.0,0) node[right] {$x$};
    \draw[axis] (0,-0.2) -- (0,2.1) node[above] {$y$};

    \coordinate (X) at (0.85,0.65);
    \coordinate (XH) at (2.25,1.25);

    \node[point] at (X) {};
    \node[point] at (XH) {};
    \draw[vector] (X) -- node[above left=-1pt] {$\h$} (XH);

    \node[below left=1pt] at (X) {$\x$};
    \node[right] at (XH) {$\x+\h$};
    \node[above] at (1.35,2.25) {domain};
  \end{scope}

  \begin{scope}[xshift=4.2cm]
    \draw[axis] (-0.25,0) -- (3.45,0) node[right] {$x$};
    \draw[axis] (0,-0.2) -- (0,2.35) node[above] {$y$};

    \coordinate (FX) at (0.65,0.72);
    \coordinate (LIN) at (2.18,1.45);
    \coordinate (ACT) at (2.98,1.95);

    \node[point] at (FX) {};
    \node[approxpoint] at (LIN) {};
    \node[point] at (ACT) {};

    \draw[linear] (FX) -- node[below right=-1pt] {$D\f(\x)\h$} (LIN);
    \draw[error] (LIN) -- node[below right=-1pt] {$\mathbf r(\h)$} (ACT);

    \node[left] at (FX) {$\f(\x)$};
    \node[above left=2pt] at (LIN) {$\f(\x)+D\f(\x)\h$};
    \node[above right=1pt] at (ACT) {$\f(\x+\h)$};
    \node[above] at (1.6,2.25) {codomain};
  \end{scope}
\end{tikzpicture}

  \caption{$D\f(\x)\h$ is the linear part; $\mathbf{r}(\h)=o(\norm{\h}_2)$ is the error.}
  \label{fig:total-derivative-linear-approximation}
\end{figure}

\begin{theorem} \label{thm:total-derivative-unique-and-continuous}\label{thm:several-variable-differentiation-main-source-todo} \hlt{Basic Consequences of Total Differentiability}\\
If $\f:E\to\R^m$ is differentiable at $\x\in E$, then $D\f(\x)$ is unique and $\f$ is continuous at $\x$.
\begin{proof}
If $A_1,A_2$ both satisfy Definition \ref{def:total-derivative-rn-rm}, set $B=A_1-A_2$. Then $\norm{B\h}_2/\norm{\h}_2\to0$ as $\h\to\ze$.\\
For fixed $\h\neq\ze$, applying this to $t\h$ and using linearity makes the quotient independent of $t$, so $B\h=\ze$.\\
Hence $B=0$. For continuity, write $\f(\x+\h)-\f(\x)=D\f(\x)\h+\mathbf{r}(\h)$.\\
The operator-norm estimate gives $\norm{D\f(\x)\h}_2\leq\normop{D\f(\x)}\norm{\h}_2$, and remainder is $o(\norm{\h}_2)$, so $\f(\x+\h)\to\f(\x)$.\\
\hlt{Structural Commentary}: The derivative is a rigid first-order linear part.\\
\hlt{Key Insight}: A remainder smaller than $\norm{\h}$ cannot hide a second linear map.
\end{proof}
\end{theorem}

\begin{example} \label{ex:derivative-of-linear-map} \hlt{Derivative of a Linear Map}\\
If $A\in\LL(\R^n,\R^m)$ and $\f(\x)=A\x$, then $D\f(\x)=A$ for every $\x\in\R^n$.\\
The remainder in Definition \ref{def:total-derivative-rn-rm} is identically zero.
\end{example}

\begin{theorem} \label{thm:chain-rule-several-variables} \hlt{Chain Rule}\\
If $\f:E\subseteq\R^n\to\R^m$ is differentiable at $\x_0$, $\g:G\subseteq\R^m\to\R^k$ is differentiable at $\f(\x_0)$, and $\f(E)\subseteq G$, then $\g\circ\f$ is differentiable at $\x_0$ and $D(\g\circ\f)(\x_0)=D\g(\f(\x_0))\,D\f(\x_0)$.
\begin{proof}
Let $A=D\f(\x_0)$, $B=D\g(\f(\x_0))$, $\mathbf{u}(\h)=\f(\x_0+\h)-\f(\x_0)-A\h$, and $\mathbf{v}(\mathbf{k})=\g(\f(\x_0)+\mathbf{k})-\g(\f(\x_0))-B\mathbf{k}$.\\
For $\mathbf{k}=\f(\x_0+\h)-\f(\x_0)$, differentiability of $\f$ gives $\mathbf{k}\to\ze$ and $(\g\circ\f)(\x_0+\h)-(\g\circ\f)(\x_0)-BA\h =B\mathbf{u}(\h)+\mathbf{v}(\mathbf{k})$, where $\norm{\mathbf{u}(\h)}_2=o(\norm{\h}_2)$ and $\norm{\mathbf{v}(\mathbf{k})}_2=o(\norm{\mathbf{k}}_2)=o(\norm{\h}_2)$. Divide by $\norm{\h}_2$.\\
\hlt{Structural Commentary}: The first-order parts compose as linear maps.\\
\hlt{Key Insight}: The output increment of $\f$ is small enough to feed into the remainder estimate for $\g$.
\end{proof}
\end{theorem}

\begin{theorem} \label{thm:several-variable-componentwise-differentiability} \hlt{Componentwise Differentiability}\\
Let $\f=(f_1,\ldots,f_m):E\subseteq\R^n\to\R^m$ and let $\pi_i:\R^m\to\R$ be the $i$th coordinate projection. Then $\f$ is differentiable at $\x\in E$ iff each $f_i$ is differentiable at $\x$, and $D f_i(\x)=\pi_i\circ D\f(\x)$ when these derivatives exist.
\begin{proof}
\begin{enumerate}[label=\roman*.]
\setlength{\itemsep}{0pt}
\item[$(\Rightarrow)$] If $\f$ is differentiable, write $\f(\x+\h)-\f(\x)=D\f(\x)\h+\mathbf r(\h)$ with $\norm{\mathbf r(\h)}_2=o(\norm{\h}_2)$.\\
Applying $\pi_i$ gives $f_i(\x+\h)-f_i(\x)=\pi_i(D\f(\x)\h)+\pi_i(\mathbf r(\h))$ and $\abs{\pi_i(\mathbf r(\h))}\leq\norm{\mathbf r(\h)}_2=o(\norm{\h}_2)$.\\
Thus $D f_i(\x)=\pi_i\circ D\f(\x)$.
\item[$(\Leftarrow)$] Suppose $D f_i(\x)=A_i\in\LL(\R^n,\R)$ for each $i$. Define $A\in\LL(\R^n,\R^m)$ by $A\h=(A_1\h,\ldots,A_m\h)$.\\
With $r_i(\h)=f_i(\x+\h)-f_i(\x)-A_i\h$, each $r_i(\h)=o(\norm{\h}_2)$.\\
So $\mathbf r(\h)=(r_1(\h),\ldots,r_m(\h))$ satisfies $\norm{\mathbf r(\h)}_2\leq\sum_i\abs{r_i(\h)}=o(\norm{\h}_2)$. Definition \ref{def:total-derivative-rn-rm} gives $D\f(\x)=A$.
\end{enumerate}
\hlt{Structural Commentary}: Component derivatives assemble into one linear map.\\
\hlt{Key Insight}: Components test the derivative; the total derivative packages the tests together.
\end{proof}
\end{theorem}

\begin{definition} \label{def:partials-gradient-directional-derivatives} \hlt{Partial, Gradient, and Directional Derivatives}\\
Let $E\subseteq\R^n$ be open, $\f=(f_1,\ldots,f_m):E\to\R^m$, and let $\mathbf{e}_1,\ldots,\mathbf{e}_n$ be the standard basis of $\R^n$.\\
If indicated limits exist, \hlt{partial derivative} is $(D_jf_i)(\x)\defeq\lim_{t\to0}\frac{f_i(\x+t\mathbf{e}_j)-f_i(\x)}{t}$ for $1\leq i\leq m,\ 1\leq j\leq n$.\\
\hlt{Directional derivative} is $(D_{\mathbf{u}}f)(\x)\defeq\lim_{t\to0}\frac{f(\x+t\mathbf{u})-f(\x)}{t}$ for $\mathbf{u}\in\R^n$; when only the direction is meant, one often assumes $\norm{\mathbf{u}}_2=1$. If $m=1$ and all partials exist, the \hlt{gradient} is $(\nabla f)(\x)\defeq \sum\limits_{j=1}^{n}(D_jf)(\x)\mathbf{e}_j$.
\end{definition}

\begin{perspective} \label{per:partials-are-coordinate-tests} \hlt{Partials Are Coordinate Tests}\\
Partials and directional derivatives test selected lines; the total derivative controls all small increments by one linear map. Theorem \ref{thm:total-derivative-matrix-partials} extracts partials, and Theorem \ref{thm:c1-iff-continuous-partials} gives a continuity-based converse.
\end{perspective}

\begin{theorem} \label{thm:total-derivative-matrix-partials} \hlt{Matrix of the Total Derivative}\\
If $\f:E\to\R^m$ is differentiable at $\x$, then every $(D_jf_i)(\x)$ exists and, in standard bases,
\begin{align}
[D\f(\x)] =
\begin{bmatrix}
(D_1f_1)(\x) & \cdots & (D_nf_1)(\x)\\
\vdots & & \vdots\\
(D_1f_m)(\x) & \cdots & (D_nf_m)(\x)
\end{bmatrix} \nonumber
\end{align}
If $m=1$ and $f$ is differentiable at $\x$, then $(D_{\mathbf{u}}f)(\x)=(\nabla f)(\x)\cdot\mathbf{u}$ for every $\mathbf{u}\in\R^n$.
\begin{proof}
Set $\h=t\mathbf{e}_j$ in Definition \ref{def:total-derivative-rn-rm}; limit is $D\f(\x)\mathbf{e}_j$, so components form column $j$ (Proposition \ref{prop:matrix-representation-action}).\\
For $m=1$, set $\h=t\mathbf{u}$, identify linear functional with dot product against $\nabla f(\x)$ (Theorem \ref{thm:riesz-representation-finite-dimensional}).\\
\hlt{Structural Commentary}: Partials are coordinate columns of $D\f(\x)$.\\
\hlt{Key Insight}: Differentiability controls all directions at once; partial derivatives record its coordinate columns.
\end{proof}
\end{theorem}

\begin{remark} \label{rem:several-variable-row-column-conventions} \hlt{Row and Column Conventions}\\
With the coordinate-column convention from Proposition \ref{prop:matrix-representation-action}, the Jacobian matrix above is the standard-basis matrix of the linear map $D\f(\x)$, and its $j$th column is $D\f(\x)\mathbf{e}_j$.
\end{remark}

\begin{theorem} \label{thm:bounded-derivative-convex-mean-value} \hlt{Mean-Value Inequality on Convex Sets}\\
Let $E\subseteq\R^n$ be open and convex.\\
If $\f:E\to\R^m$ is differentiable and $\normop{D\f}\leq M$ on $E$, then for all $\a,\b\in E$, $\norm{\f(\b)-\f(\a)}_2\leq M\norm{\b-\a}_2$.\\
If $D\f=0$ on $E$, then $\f$ is constant.
\begin{proof}
Fix $\a,\b\in E$, set $\gamma(t)=(1-t)\a+t\b$, and let $\mathbf{g}=\f\circ\gamma$. Then $\gamma([0,1])\subseteq E$ and $\norm{\mathbf{g}'(t)}_2=\norm{D\f(\gamma(t))(\b-\a)}_2\leq M\norm{\b-\a}_2$. Apply Theorem \ref{thm:vector-valued-mean-value-inequality}; the constant case is $M=0$.\\
\hlt{Structural Commentary}: Restricting to a line segment reduces the estimate to one variable.\\
\hlt{Key Insight}: A uniform bound on the derivative bounds every secant displacement.
\end{proof}
\end{theorem}

\begin{definition} \label{def:c1-map} \hlt{$C^1$ Mapping}\\
Let $E\subseteq\R^n$ be open. A differentiable map $\f:E\to\R^m$ is \hlt{continuously differentiable}, or $C^1$, if $\x\mapsto D\f(\x)$ is continuous from $E$ into $\LL(\R^n,\R^m)$ in operator norm.
\end{definition}

\begin{perspective} \label{per:c1-operator-norm-continuity} \hlt{$C^1$ Uses the Operator-Norm Metric}\\
$C^1$ continuity measures $\x\mapsto D\f(\x)$ in operator norm. In coordinates this is spectral-norm continuity, equivalent here to continuity of the derivative matrix entries.
\end{perspective}

\begin{figure}[H]
  \centering
  \begin{tikzpicture}[
  scale=0.95,
  >=Stealth,
  every node/.style={font=\scriptsize},
  segment/.style={->, line width=0.85pt, blue!70!black},
  direct/.style={->, line width=0.55pt, black!40, dashed},
  point/.style={circle, fill=black, inner sep=1.15pt}
]
  \coordinate (p0) at (0,0);
  \coordinate (p1) at (1.55,0);
  \coordinate (p2) at (2.55,0.72);
  \coordinate (p3) at (2.55,2.02);

  \draw[direct] (p0) -- node[above left=2pt] {$\h$} (p3);
  \draw[segment] (p0) -- node[below=2pt] {$h_1\mathbf{e}_1$} (p1);
  \draw[segment] (p1) -- node[below right=1pt] {$h_2\mathbf{e}_2$} (p2);
  \draw[segment] (p2) -- node[right=2pt] {$h_3\mathbf{e}_3$} (p3);

  \node[point] at (p0) {};
  \node[point] at (p1) {};
  \node[point] at (p2) {};
  \node[point] at (p3) {};

  \node[left=3pt] at (p0) {$\x^{(0)}=\x$};
  \node[below=5pt] at (p1) {$\x^{(1)}$};
  \node[right=4pt] at (p2) {$\x^{(2)}$};
  \node[above=4pt] at (p3) {$\x^{(3)}=\x+\h$};
\end{tikzpicture}

  \caption{Coordinate increments telescope $\x$ to $\x+\h$.}
  \label{fig:coordinate-polygon-partial-path}
\end{figure}

\begin{theorem} \label{thm:c1-iff-continuous-partials} \hlt{$C^1$ Means Continuous Partials}\\
A map $\f:E\subseteq\R^n\to\R^m$ is $C^1$ iff all $D_jf_i$ exist and are continuous on $E$.
\begin{proof}
If $\f$ is $C^1$, Theorem \ref{thm:total-derivative-matrix-partials} makes each $D_jf_i$ a continuous matrix entry of $D\f$.\\
Conversely reduce componentwise to $m=1$. Fix $\x\in E$ and take $\h$ small enough that the coordinate polygon from $\x$ to $\x+\h$ lies in $E$. Put $\x^{(0)}=\x$ and $\x^{(j)}=\x+h_1\mathbf{e}_1+\cdots+h_j\mathbf{e}_j$. By one-variable MVT on each coordinate segment, $f(\x+\h)-f(\x)=\sum\limits_{j=1}^{n}h_j(D_jf)(\xi_j)$ for points $\xi_j$ on corresponding segments. Continuity of partials at $\x$ gives $\abs{f(\x+\h)-f(\x)-\sum\limits_{j=1}^{n}h_j(D_jf)(\x)} \leq \sum\limits_{j=1}^{n}\abs{h_j}\abs{(D_jf)(\xi_j)-(D_jf)(\x)} =o(\norm{\h}_2)$.\\
Hence $Df(\x)$ has row matrix $[(D_1f)(\x)\ \cdots\ (D_nf)(\x)]$. Since operator norm is controlled by the Euclidean norm of the matrix entries, continuity of the partials gives continuity of $\x\mapsto Df(\x)$.\\
\hlt{Structural Commentary}: Continuous partial derivatives assemble into a continuous total derivative.\\
\hlt{Key Insight}: Coordinate-by-coordinate increments work only because the partials vary continuously.
\end{proof}
\end{theorem}

Higher derivatives repeat the same idea: once $D\f$ is itself a map into the finite-dimensional vector space $\LL(\R^n,\R^m)$, we differentiate it again. The linear-algebra viewpoint matters here because the second derivative is naturally bilinear, not merely an array of second partial derivatives.

\begin{definition} \label{def:higher-order-partials-and-c2} \hlt{Higher Derivatives and $C^r$}\\
Let $E\subseteq\R^n$ be open, $\f:E\to\R^m$, $\LL^r(\R^n,\R^m)$ be space of $r$-linear maps $(\R^n)^r\to\R^m$. Set $D^1\f\defeq D\f$. Recursively, if $D^{r-1}\f:E\to\LL^{r-1}(\R^n,\R^m)$ differentiable at $\x$, let $D^r\f(\x)(\mathbf{v}_1,\ldots,\mathbf{v}_r) \defeq \bigl(D(D^{r-1}\f)(\x)\mathbf{v}_1\bigr)(\mathbf{v}_2,\ldots,\mathbf{v}_r)$.\\
For real-valued $f$, write $D_{ij}f\defeq D_i(D_jf)$ when defined. The map $\f$ is $C^r$ if derivatives through order $r$ exist and $D^r\f:E\to\LL^r(\R^n,\R^m)$ is continuous; it is $C^\infty$ if it is $C^r$ for all $r$.
\end{definition}

\begin{theorem} \label{thm:second-derivative-partials-and-continuity} \hlt{Second Derivatives and Second Partials}\\
Let $\f=(f_1,\ldots,f_m):E\subseteq\R^n\to\R^m$ and $p\in E$.
\begin{enumerate}[label=\roman*.]
\setlength{\itemsep}{0pt}
\item[$(\Rightarrow)$] If $D^2\f(p)$ exists, then each $D_{ij}f_\ell(p)$ exists and $D^2 f_\ell(p)(\mathbf{e}_i,\mathbf{e}_j)=D_{ij}f_\ell(p)$.
\item[$(\Leftarrow)$] If all $D_{ij}f_\ell$ exist near $p$ and are continuous at $p$, then $D^2\f(p)$ exists.
\end{enumerate}
\begin{proof}
\begin{enumerate}[label=\roman*.]
\setlength{\itemsep}{0pt}
\item[$(\Rightarrow)$] Apply Theorem \ref{thm:several-variable-componentwise-differentiability} to $D\f:E\to\LL(\R^n,\R^m)$. The entries of $D\f$ are $D_jf_\ell$ by Theorem \ref{thm:total-derivative-matrix-partials}; differentiating the $(\ell,j)$ entry in the $\mathbf{e}_i$ direction gives $D_{ij}f_\ell(p)$ and the displayed formula.
\item[$(\Leftarrow)$] Matrix entries of $D\f$ are functions $D_jf_\ell$. Since first partials $D_i(D_jf_\ell)=D_{ij}f_\ell$ exist near $p$ and are continuous at $p$, coordinate criterion at $p$ gives differentiability of $D\f$ there. Hence $D^2\f(p)$ exists.
\end{enumerate}
\hlt{Structural Commentary}: Second partials are matrix entries of the derivative of the first-derivative matrix.\\
\hlt{Key Insight}: Continuous second partials make the first-derivative matrix differentiable.
\end{proof}
\end{theorem}

\begin{figure}[H]
  \centering
  \begin{tikzpicture}[
  scale=0.95,
  >=Stealth,
  every node/.style={font=\scriptsize},
  edge/.style={black!70, line width=0.7pt},
  stepvec/.style={->, blue!70!black, line width=0.75pt},
  point/.style={circle, fill=black, inner sep=1.2pt},
  sign/.style={font=\normalsize}
]
  \coordinate (p) at (0,0);
  \coordinate (v) at (4.1,0.35);
  \coordinate (w) at (1.15,2.25);
  \coordinate (vw) at (5.25,2.60);

  \fill[blue!4] (p) -- (v) -- (vw) -- (w) -- cycle;
  \draw[edge] (p) -- (v) -- (vw) -- (w) -- cycle;

  \node[point] at (p) {};
  \node[point] at (v) {};
  \node[point] at (w) {};
  \node[point] at (vw) {};

  \node[below left=2pt] at (p) {$p$};
  \node[below=4pt] at (v) {$p+t\mathbf{v}$};
  \node[above left=2pt] at (w) {$p+t\mathbf{w}$};
  \node[above right=2pt] at (vw) {$p+t\mathbf{v}+t\mathbf{w}$};

  \node[sign] at (0.62,0.38) {$+$};
  \node[sign] at (3.65,0.75) {$-$};
  \node[sign] at (1.42,1.90) {$-$};
  \node[sign] at (4.72,2.25) {$+$};

  \draw[stepvec] (0.45,-0.40) -- node[below=2pt] {$t\mathbf{v}$} (3.65,-0.13);
  \draw[stepvec] (-0.55,0.10) -- node[left=2pt] {$t\mathbf{w}$} (0.60,2.35);
\end{tikzpicture}

  \caption{The signed vertex sum on a small parallelogram is symmetric in $\mathbf{v}$ and $\mathbf{w}$.}
  \label{fig:signed-parallelogram-second-derivative}
\end{figure}

\begin{theorem} \label{thm:second-derivative-symmetric} \hlt{Symmetry of the Second Derivative}\\
If $D^2\f(p)$ exists, then $D^2\f(p)(\mathbf{v},\mathbf{w})=D^2\f(p)(\mathbf{w},\mathbf{v})$ for all $\mathbf{v},\mathbf{w}\in\R^n$.
\begin{proof}
Reduce to scalar $f$ by Theorem \ref{thm:several-variable-componentwise-differentiability}.\\
For small $t$, set $\Delta(t,\mathbf{v},\mathbf{w})\defeq f(p+t\mathbf{v}+t\mathbf{w})-f(p+t\mathbf{v})-f(p+t\mathbf{w})+f(p)$.\\
Then $\Delta(t,\mathbf{v},\mathbf{w})=\Delta(t,\mathbf{w},\mathbf{v})$; see Figure \ref{fig:signed-parallelogram-second-derivative}. The one-variable mean value theorem and differentiability of $Df$ at $p$ give $D^2 f(p)(\mathbf{v},\mathbf{w}) =\lim\limits_{t\to0}\frac{\Delta(t,\mathbf{v},\mathbf{w})}{t^2}$. The same formula with $\mathbf{v},\mathbf{w}$ reversed proves symmetry.\\
\hlt{Structural Commentary}: A pointwise second derivative is already forced to be a symmetric bilinear form.\\
\hlt{Key Insight}: The parallelogram difference is symmetric in the two directions.
\end{proof}
\end{theorem}

\begin{corollary} \label{cor:mixed-partials-from-second-derivative} \hlt{Mixed Partials from $D^2$}\\
If $D^2\f(p)$ exists, then $D_{ij}f_\ell(p)=D_{ji}f_\ell(p)$ for all components and all $i,j$.
\begin{proof}
By Theorem \ref{thm:second-derivative-partials-and-continuity}, $D_{ij}f_\ell(p)=D^2f_\ell(p)(\mathbf{e}_i,\mathbf{e}_j)$; and $D_{ji}f_\ell(p)=D^2f_\ell(p)(\mathbf{e}_j,\mathbf{e}_i)$.\\
Theorem \ref{thm:second-derivative-symmetric} makes the two right-hand sides equal.\\
\hlt{Structural Commentary}: Mixed partials are coordinate values of one bilinear map.\\
\hlt{Key Insight}: Symmetry of $D^2f_\ell(p)$ gives equality of mixed entries.
\end{proof}
\end{corollary}

\begin{corollary} \label{cor:higher-derivatives-symmetric} \hlt{Symmetry of Higher Derivatives}\\
If $D^r\f(p)$ exists, then the $r$-linear map $D^r\f(p)$ is unchanged by permuting its arguments.
\begin{proof}
It suffices to prove invariance under adjacent transpositions. Fix all arguments except positions $k,k+1$.\\
Resulting two-variable second derivative is $(\mathbf{u},\mathbf{v})\longmapsto D^r\f(p)(\mathbf{v}_1,\ldots,\mathbf{v}_{k-1},\mathbf{u},\mathbf{v},\mathbf{v}_{k+2},\ldots,\mathbf{v}_r)$, so Theorem \ref{thm:second-derivative-symmetric} allows $\mathbf{u}$ and $\mathbf{v}$ to be swapped. Since adjacent transpositions generate all permutations, $D^r\f(p)$ is symmetric.\\
\hlt{Structural Commentary}: Higher derivatives are iterated second-derivative symmetries.\\
\hlt{Key Insight}: Adjacent swaps generate every permutation.
\end{proof}
\end{corollary}

The next rectangle argument is the classical coordinate route to equality of mixed partials. It is weaker in hypothesis than global $C^2$, but it proves the same compatibility of coordinate second derivatives.

\begin{theorem} \label{thm:rectangle-mean-value-mixed-partial} \hlt{Rectangle Mean Value Formula}\\
Let $E\subseteq\R^2$ be open, and suppose $D_1f$ and $D_{21}f$ exist on $E$. For a closed coordinate rectangle $Q\subseteq E$ with opposite vertices $(a,b)$ and $(a+h,b+k)$, $h,k\neq0$, set $\Delta(f,Q)\defeq f(a+h,b+k)-f(a+h,b)-f(a,b+k)+f(a,b)$.\\
Then some $(x,y)\in\operatorname{int}Q$ satisfies $\Delta(f,Q)=hk(D_{21}f)(x,y)$.
\begin{proof}
Set $u(t)=f(t,b+k)-f(t,b)$.\\
By Corollary \ref{cor:mean-value-theorem}, for some $x \in (a, a+h)$, $\Delta(f,Q)=u(a+h)-u(a)=h\bigl[(D_1f)(x,b+k)-(D_1f)(x,b)\bigr]$.\\
Apply same corollary to $s\mapsto(D_1f)(x,s)$; for some $y \in (b, b+k)$, the bracket equals $k(D_{21}f)(x,y)$.\\
Hence $(x,y)\in\operatorname{int}Q$ and the formula follows.\\
\hlt{Structural Commentary}: The rectangle difference is a two-step one-variable mean value calculation.\\
\hlt{Key Insight}: Mixed second derivatives measure the first-order change of a first partial across the other coordinate.
\end{proof}
\end{theorem}

\begin{theorem} \label{thm:equality-of-mixed-partials} \hlt{Equality of Mixed Partials}\\
Let $(a,b)$ lie in an open set $E\subseteq\R^2$. If $D_1f$, $D_2f$, and $D_{21}f$ exist on $E$, and $D_{21}f$ is continuous at $(a,b)$, then $D_{12}f(a,b)$ exists and $(D_{12}f)(a,b)=(D_{21}f)(a,b)$. Hence $D_{12}f=D_{21}f$ when $f\in C^2(E)$.
\begin{proof}
Let $A=(D_{21}f)(a,b)$.\\
Given $\varepsilon>0$, continuity of $D_{21}f$ and Theorem \ref{thm:rectangle-mean-value-mixed-partial} give, for all small $h,k\neq0$, $\abs{\frac{\Delta(f,Q)}{hk}-A}<\varepsilon$.\\
Fix such an $h$ and let $k\to0$. Since $D_2f$ exists at $(a+h,b)$ and $(a,b)$, $\frac{\Delta(f,Q)}{hk}\to \frac{(D_2f)(a+h,b)-(D_2f)(a,b)}{h}$.\\
Thus this last quotient differs from $A$ by at most $\varepsilon$ for all sufficiently small $h\neq0$.\\
Since $\varepsilon$ was arbitrary, the quotient tends to $A$ as $h\to0$, so $D_{12}f(a,b)$ exists and equals $A$.\\
\hlt{Structural Commentary}: Continuity of $D_{21}f$ controls local rectangle quotients.\\
\hlt{Key Insight}: The missing mixed partial is recovered by taking the two one-variable limits in the reverse order.
\end{proof}
\end{theorem}

\begin{perspective} \label{per:mixed-partials-need-hypotheses} \hlt{Mixed Partials Need Hypotheses}\\
$D_{12}f$ and $D_{21}f$ are different iterated limits. Theorems \ref{thm:second-derivative-symmetric} and \ref{thm:equality-of-mixed-partials} give hypotheses for equality.
\end{perspective}

\begin{theorem} \label{thm:differentiation-under-riemann-stieltjes-integral} \hlt{Differentiation Under a Riemann--Stieltjes Integral}\\
Let $\alpha$ be increasing on $[a,b]$, let $\varphi:[a,b]\times[c,d]\to\R$, and put $\varphi^t(x)\defeq\varphi(x,t)$.\\
Assume $\varphi^t\in\mathcal{R}(\alpha)$ for all $t$ and $D_2\varphi(\cdot,t)\rightrightarrows D_2\varphi(\cdot,s)$ on $[a,b]$ as $t\to s\in(c,d)$.\\
If $F(t)\defeq\int_a^b \varphi(x,t)\,d\alpha(x)$, then $(D_2\varphi)^s\in\mathcal{R}(\alpha)$ and $F'(s)=\int_a^b (D_2\varphi)(x,s)\,d\alpha(x)$.
\begin{proof}
For $t\neq s$, set $\psi_t=(\varphi^t-\varphi^s)/(t-s)$. Corollary \ref{cor:mean-value-theorem} and the uniform hypothesis give $\psi_t\rightrightarrows(D_2\varphi)^s$.\\
Since $\psi_t\in\mathcal{R}(\alpha)$, Theorem \ref{thm:uniform-convergence-and-riemann-stieltjes-integration} gives $\frac{F(t)-F(s)}{t-s}=\int_a^b\psi_t\,d\alpha \longrightarrow \int_a^b(D_2\varphi)(x,s)\,d\alpha(x)$.\\
\hlt{Structural Commentary}: Uniform parameter control lets limits pass under the integral.\\
\hlt{Key Insight}: The theorem is exactly the commutation of a derivative limit with a Riemann--Stieltjes integral.
\end{proof}
\end{theorem}

\begin{example} \label{ex:gaussian-fourier-integral-differentiation} \hlt{A Gaussian Fourier Integral}\\
Define $f(t)=\int_{-\infty}^{\infty}e^{-x^2}\cos(xt)\,dx$ and $g(t)=-\int_{-\infty}^{\infty}xe^{-x^2}\sin(xt)\,dx$.\\
Then $f'(t)=g(t)$, and integration by parts gives $2f'(t)+tf(t)=0$. Since $f(0)=\sqrt{\pi}$, $f(t)=\sqrt{\pi}\exp\!\left(-\frac{t^2}{4}\right)$.
\end{example}

\begin{definition} \label{def:jacobian-determinant} \hlt{Jacobian}\\
If $\f:E\subseteq\R^n\to\R^n$ is differentiable at $\x$, define $J_{\f}(\x)\defeq \det D\f(\x) =\frac{\partial(y_1,\ldots,y_n)}{\partial(x_1,\ldots,x_n)}$ when $f=(y_1,\ldots,y_n)$.\\
By Theorem \ref{thm:determinant-basic-laws}, $D\f(\x)$ is invertible if and only if $J_{\f}(\x)\neq0$. The absolute value $\abs{J_{\f}(\x)}$ is the local version of determinant volume scaling from Theorem \ref{thm:linear-map-volume-scaling-determinant}, used again in the change-of-variables theorem.
\end{definition}
